\section{Introduction}

Expressions that mix gradients with respect to position and momentum can suggest a connection between phase-space geometry and quantum noncommutativity. A representative trial expression is
\begin{equation}
\nabla_{\bm{x}}\times\nabla_{\bm{p}}
=
\frac{\hbar^2}{2\pi}.
\label{eq:bad}
\end{equation}
The intuition behind Eq.~\eqref{eq:bad} is understandable: Planck's constant controls the relation between conjugate observables, and antisymmetric operations such as commutators, wedge products, and curls all encode some form of order dependence or circulation. These operations are nevertheless mathematically different. Substituting one for another without specifying domains, tensor type, dimensions, and operator ordering produces false identities.

The purpose of this note is therefore diagnostic rather than legislative. We do not propose a universal ``normalization'' rule for powers of \(\hbar\). Instead, we establish a hierarchy of correct statements and identify the assumptions under which each one applies. The central distinctions are:
\begin{enumerate}[label=(\roman*)]
  \item mixed coordinate derivatives commute;
  \item a formal cross product of two gradient operators need not vanish;
  \item the intrinsic antisymmetric structure on canonical phase space is a Poisson bivector;
  \item quantum noncommutativity is expressed by an operator or star commutator;
  \item quantized vorticity is a topological statement in physical space, not a canonical phase-space identity.
\end{enumerate}
